% =====================================================================================
% B-transport.tex — PHẦN II của Nguyễn Minh Toàn: tầng giao vận UDP (Transport)
% =====================================================================================
\PhanLon{II. NỘI DUNG CÁ NHÂN}
\textit{Phân hệ tầng giao vận UDP (Transport).}

\section{Biểu đồ và giao diện nội dung cá nhân thực hiện}
Phân hệ Transport xây dựng \textbf{tầng giao vận tin cậy trên UDP từ con số không}, chỉ với \code{Socket} kiểu
\code{Dgram}: đánh số thứ tự, ack bằng bitfield, gửi lại, 4 kênh độc lập, phân mảnh, ước lượng RTT, kiểm soát
nghẽn, chống giả mạo và bộ giả lập mạng. Tầng này chỉ biết mảng byte --- không biết nội dung là snapshot hay
chat --- và không tham chiếu \code{UnityEngine}.

% -------------------------------------------------------------------------------------
\subsection{Kiến trúc thư viện Transport}
\begin{figure}[H]
\centering
\begin{tikzpicture}[x=1cm, y=1cm, lay/.style={box, text width=3.4cm, minimum height=1.3cm},
                    rowname/.style={font=\small\bfseries, align=left, text width=2.5cm, anchor=west}]
  \node[rowname] at (0,0) {Ứng dụng};
  \node[lay, fill=cOrange] (u1) at (4.6,0)  {\textbf{Client Unity}\\{\footnotesize NetClientBootstrap}};
  \node[lay, fill=cOrange] (u2) at (8.4,0)  {\textbf{Game server}\\{\footnotesize NetServerBootstrap}};
  \node[lay, fill=cOrange] (u3) at (12.2,0) {\textbf{LoadTest, xUnit}\\{\footnotesize bot và kiểm thử}};
  \node[font=\footnotesize] at (8.4,-1.05) {\texttt{ITransportClient} / \texttt{ITransportServer} -- chỉ trao đổi byte và id kênh};
  \node[rowname] at (0,-2.1) {Mặt tiền API};
  \node[lay, fill=cBlue]   (a1) at (4.6,-2.1)  {\textbf{UdpTransportClient}\\{\footnotesize kết nối, gửi lại yêu cầu}};
  \node[lay, fill=cBlue]   (a2) at (8.4,-2.1)  {\textbf{UdpTransportServer}\\{\footnotesize bảng kết nối, RateLimiter}};
  \node[lay, fill=cPurple] (a3) at (12.2,-2.1) {\textbf{LoopbackTransport}\\{\footnotesize trong bộ nhớ, cho test}};
  \node[rowname] at (0,-3.8) {Mỗi kết nối};
  \node[lay, fill=cBlue] (k1) at (4.6,-3.8)  {\textbf{Connection}\\{\footnotesize trạng thái, RTT, keep-alive}};
  \node[lay, fill=cBlue] (k2) at (8.4,-3.8)  {\textbf{ReliabilityLayer}\\{\footnotesize seq, ack, bitfield, gửi lại}};
  \node[lay, fill=cBlue] (k3) at (12.2,-3.8) {\textbf{ChannelSet}\\{\footnotesize 4 kênh, sequence riêng}};
  \node[lay, fill=cBlue] (k4) at (4.6,-5.3)  {\textbf{FragmentAssembler}\\{\footnotesize ghép mảnh, 8 nhóm, 2 s}};
  \node[lay, fill=cBlue] (k5) at (8.4,-5.3)  {\textbf{CongestionControl}\\{\footnotesize GOOD / BAD}};
  \node[lay, fill=cBlue] (k6) at (12.2,-5.3) {\textbf{FlowControl}\\{\footnotesize tối đa 64 gói chưa ack}};
  \node[rowname] at (0,-7.1) {Socket};
  \node[lay, fill=cGreen] (s1) at (4.6,-7.1)  {\textbf{UdpPeer}\\{\footnotesize socket không chặn, 1024 gói/Poll}};
  \node[lay, fill=cGray]  (s2) at (8.4,-7.1)  {\textbf{BufferPool, PacketLogger}\\{\footnotesize không cấp phát, ghi .ifpcap}};
  \node[lay, fill=cGray]  (s3) at (12.2,-7.1) {\textbf{NetworkSimulator}\\{\footnotesize mất gói, trễ, jitter}};
  \node[font=\footnotesize] at (8.4,-8.25) {\texttt{System.Net.Sockets.Socket} (\texttt{SocketType.Dgram}) -- hệ điều hành};
\end{tikzpicture}
\caption{Biểu đồ thành phần thư viện Ironfront.Net.Transport}
\label{fig:B-thanhphan}
\end{figure}

% -------------------------------------------------------------------------------------
\subsection{Định dạng gói tin GSP}
\begin{figure}[H]
\centering
\begin{tikzpicture}[x=1cm, y=1cm, fld/.style={draw=cLine, minimum height=1.2cm, align=center, font=\scriptsize, inner sep=1pt, anchor=west}]
  % 16 byte = 16 cm... dùng 0.98 cm / byte
  \def\b{0.98}
  \node[fld, fill=cBlue,   minimum width=2*\b cm] (p)  at (0,0)      {\textbf{protocolId}\\u16 0x4946};
  \node[fld, fill=cBlue,   minimum width=1*\b cm] (t)  at (p.east)   {\textbf{type}\\u8};
  \node[fld, fill=cBlue,   minimum width=1*\b cm] (f)  at (t.east)   {\textbf{flags}\\u8};
  \node[fld, fill=cGreen,  minimum width=2*\b cm] (s)  at (f.east)   {\textbf{sequence}\\u16};
  \node[fld, fill=cGreen,  minimum width=2*\b cm] (a)  at (s.east)   {\textbf{ack}\\u16};
  \node[fld, fill=cGreen,  minimum width=4*\b cm] (ab) at (a.east)   {\textbf{ackBitfield}\\u32: 32 gói trước ack};
  \node[fld, fill=cOrange, minimum width=2*\b cm] (c)  at (ab.east)  {\textbf{connection}\\\textbf{Id} u16};
  \node[fld, fill=cOrange, minimum width=2*\b cm] (l)  at (c.east)   {\textbf{payload}\\\textbf{Length}\\u16 $\le$ 1184};
  \foreach \n/\o in {p/0, t/2, f/3, s/4, a/6, ab/8, c/12, l/14} { \node[font=\scriptsize, anchor=north west] at (\n.south west) {\o}; }
  \node[font=\scriptsize, anchor=north east] at (l.south east) {16};
  % envelope + fragment header
  \node[fld, fill=cPurple, minimum width=2.4cm] (e1) at (0,-2.4)  {\textbf{channelId}\\u8: 0..3};
  \node[fld, fill=cPurple, minimum width=4.2cm] (e2) at (e1.east) {\textbf{channelSequence}\\u16, riêng từng kênh};
  \node[font=\footnotesize, anchor=north] at ($(e1.south west)!0.5!(e2.south east)+(0,-0.1)$) {Envelope kênh 3 byte (mọi gói PAYLOAD)};
  \node[fld, fill=cPurple, minimum width=3.6cm] (g1) at (8.1,-2.4)  {\textbf{fragmentGroupId}\\u16};
  \node[fld, fill=cPurple, minimum width=2.0cm] (g2) at (g1.east)   {\textbf{index}\\u8};
  \node[fld, fill=cPurple, minimum width=2.0cm] (g3) at (g2.east)   {\textbf{count}\\u8 $\le$ 64};
  \node[font=\footnotesize, anchor=north] at ($(g1.south west)!0.5!(g3.south east)+(0,-0.1)$) {Header phân mảnh 4 byte (khi cờ FRAGMENTED)};
\end{tikzpicture}
\caption{Header GSP 16 byte (little-endian), envelope kênh và header phân mảnh}
\label{fig:B-header}
\end{figure}

\begin{center}\small
\begin{tabularx}{\linewidth}{|L{2.6cm}|Y|}
\hline
\hd{Trường} & \hd{Giá trị} \\ \hline
Bit \code{flags} & 0 RELIABLE $\cdot$ 1 FRAGMENTED $\cdot$ 2 ORDERED $\cdot$ 3 COMPRESSED (dự phòng) $\cdot$ 4--7 phải bằng 0 \\ \hline
\code{packetType} & 0x01 CONNECT\_REQUEST $\cdot$ 0x02 CONNECT\_CHALLENGE $\cdot$ 0x03 CONNECT\_RESPONSE $\cdot$ 0x04 CONNECT\_ACCEPTED $\cdot$ 0x05 CONNECT\_DENIED $\cdot$ 0x06 DISCONNECT $\cdot$ 0x07 KEEPALIVE $\cdot$ 0x10 PAYLOAD $\cdot$ 0x11 FRAGMENT \\ \hline
\end{tabularx}
\end{center}

% -------------------------------------------------------------------------------------
\subsection{Bắt tay không lưu trạng thái}
\begin{figure}[H]
\centering
\begin{tikzpicture}[x=1cm, y=1cm]
  \seqactor{C}{2.2}{Client (UdpTransportClient)}
  \seqactor{S}{12.8}{Server (UdpTransportServer)}
  \seqbegin
  \seqmsg{C}{S}{CONNECT\_REQUEST \{version, joinTicket 64 byte, clientSalt u64\}}
  \seqnote{S}{S}{3}{RateLimiter theo IP; kiểm version và chỗ trống; \textbf{không lưu gì}, chưa kiểm HMAC}
  \seqret{S}{C}{CONNECT\_CHALLENGE \{serverSalt\}}
  \seqnote{C}{S}{2}{serverSalt = HMAC(serverKey, địa chỉ, cổng, clientSalt, epoch 30 s)\\client tính challengeResponse = clientSalt XOR serverSalt}
  \seqmsg{C}{S}{CONNECT\_RESPONSE \{response, clientSalt, joinTicket\}}
  \seqnote{S}{S}{3}{Tính lại salt: địa chỉ đã được chứng minh; kiểm vé (HMAC, hạn), gắn playerId, cấp connectionId}
  \seqfragbegin
  \seqret{S}{C}{CONNECT\_ACCEPTED \{connectionId, serverTick, mapId, playerId\}}
  \seqfragelse{C}{S}{vé sai / hết hạn / playerId đang chơi}
  \seqret{S}{C}{CONNECT\_DENIED \{lý do 1..6\}}
  \seqfragend{C}{S}{alt}{vé hợp lệ}
  \seqend{C,S}
\end{tikzpicture}
\caption{Bắt tay challenge -- response dùng cookie HMAC: server không cấp tài nguyên cho địa chỉ chưa được chứng minh}
\label{fig:B-battay}
\end{figure}

% -------------------------------------------------------------------------------------
\subsection{Máy trạng thái kết nối}
\begin{figure}[H]
\centering
\begin{tikzpicture}[x=1cm, y=1cm]
  \node[init] (i) at (-0.6,0) {};
  \node[st, text width=2.4cm] (d)  at (1.8,0)    {Disconnected};
  \node[st, text width=2.4cm] (c)  at (7.2,0)    {Connecting};
  \node[st, text width=2.4cm] (h)  at (12.6,0)   {Challenged};
  \node[st, text width=2.4cm, fill=cGreen] (ok) at (12.6,-3.4) {Connected};
  \draw[arr] (i) -- (d);
  \draw[arr] (d) -- node[lbl, above=3pt]{Connect()} (c);
  \draw[arr] (c) -- node[lbl, above=3pt]{nhận CHALLENGE} (h);
  \draw[arr] (h) -- node[lbl, right=3pt]{nhận CONNECT\_ACCEPTED} (ok);
  \draw[arr] (c.north) -- ++(0,0.9) -| node[lbl, pos=0.25, above=2pt]{5 s hết giờ / CONNECT\_DENIED} ($(d.north)+(0.4,0)$);
  \draw[arr] (h.north) -- ++(0,1.8) -| node[lbl, pos=0.25, above=2pt]{5 s hết giờ} ($(d.north)+(-0.4,0)$);
  \draw[arr] (ok.west) -| node[lbl, pos=0.3, below=3pt]{DISCONNECT (gửi 3 lần) / 10 s im lặng / hết lượt gửi lại} (d.south);
  \draw[arr] (ok.south east) .. controls +(1.2,-0.8) and +(1.3,0.4) .. node[lbl, right=2pt]{KEEPALIVE\\mỗi 1 s} (ok.east);
  \node[font=\footnotesize, align=left, anchor=north west] at (3.3,-0.55) {CONNECT\_REQUEST gửi lại mỗi 250 ms,\\tối đa 20 lần};
\end{tikzpicture}
\caption{Máy trạng thái \texttt{ConnectionState} của một kết nối}
\label{fig:B-trangthai}
\end{figure}

% -------------------------------------------------------------------------------------
\subsection{Cơ chế ack bằng bitfield và gửi lại}
\begin{figure}[H]
\centering
\begin{tikzpicture}[x=1cm, y=1cm, cell/.style={draw=cLine, minimum width=2.1cm, minimum height=1.1cm, align=center, font=\footnotesize, anchor=west}]
  \node[cell, fill=cGreen]  (c0) at (0,0)     {\textbf{seq 103}\\đã nhận};
  \node[cell, fill=cOrange] (c1) at (c0.east) {\textbf{seq 102}\\MẤT};
  \node[cell, fill=cGreen]  (c2) at (c1.east) {\textbf{seq 101}\\đã nhận};
  \node[cell, fill=cOrange] (c3) at (c2.east) {\textbf{seq 100}\\MẤT};
  \node[cell, fill=cGreen]  (c4) at (c3.east) {\textbf{seq 99}\\đã nhận};
  \node[cell, fill=cGreen]  (c5) at (c4.east) {\textbf{seq 98}\\đã nhận};
  \node[cell, fill=cGreen]  (c6) at (c5.east) {\textbf{seq 97}\\đã nhận};
  \foreach \n/\t in {c0/{ack = 103}, c1/{bit0 = 0}, c2/{bit1 = 1}, c3/{bit2 = 0}, c4/{bit3 = 1}, c5/{bit4 = 1}, c6/{bit5 = 1}}
    { \node[font=\footnotesize] at ($(\n.south)+(0,-0.3)$) {\t}; }
\end{tikzpicture}
\caption{Ví dụ ack bitfield: \texttt{ack = 103}, \texttt{ackBitfield = ...111010}; bên gửi biết ngay gói 102 và 100 bị mất.
Mỗi gói mang 33 thông tin ack nên không cần gói ACK riêng}
\label{fig:B-bitfield}
\end{figure}

\begin{figure}[H]
\centering
\begin{tikzpicture}[x=1cm, y=1cm]
  % ---- cột gửi
  \node[font=\small\bfseries] at (4,1.0) {Gửi};
  \node[bxB, text width=3.6cm] (a)  at (4,0)    {\texttt{Send(channel, payload, reliable)}};
  \node[dec, text width=1.7cm] (d1) at (4,-1.8) {reliable?};
  \node[dec, text width=2.2cm] (d2) at (4,-4.0) {chưa ack\\< 64 gói?};
  \node[bxO, text width=1.7cm] (q)  at (7.9,-4.0) {Tạm hoãn gói tin cậy mới};
  \node[bxB, text width=3.6cm] (k)  at (4,-6.1)  {Sao chép vào BufferPool, lưu theo sequence};
  \node[bxG, text width=3.6cm] (w)  at (4,-7.9)  {Ghi header + envelope\\gửi datagram};
  \draw[arr] (a) -- (d1);
  \draw[arr] (d1) -- node[lbl, right=2pt]{có} (d2);
  \draw[arr] (d2) -- node[lbl, above=2pt]{không} (q);
  \draw[arr] (d2) -- node[lbl, right=2pt]{có} (k);
  \draw[arr] (k) -- (w);
  \draw[arr] (d1.west) -- node[lbl, above=2pt]{không} (0.6,-1.8) |- (w.west);
  % ---- cột nhận ack / gửi lại
  \draw[cLine, dashed] (9.0,1.2) -- (9.0,-8.6);
  \node[font=\small\bfseries] at (12.6,1.0) {Nhận ack và gửi lại};
  \node[bxB, text width=3.8cm] (r)  at (12.6,0)    {Gói đến mang ack + bitfield};
  \node[bxB, text width=3.8cm] (x)  at (12.6,-1.8) {Đánh dấu đã nhận, trả buffer; mẫu RTT chỉ từ lần gửi đầu (Karn)};
  \node[bxB, text width=3.8cm] (t)  at (12.6,-4.0) {\texttt{Poll()}: gói quá RTO\\(30--1000 ms)};
  \node[dec, text width=2.2cm] (d3) at (12.6,-6.1) {đã gửi lại\\> 10 lần?};
  \node[bxO, text width=3.8cm] (z)  at (12.6,-7.9) {Đóng kết nối\\\texttt{TransportError}};
  \draw[arr] (r) -- (x);
  \draw[arr] (t) -- (d3);
  \draw[arr] (d3) -- node[lbl, right=2pt]{có} (z);
  \draw[arr] (d3.west) -- node[lbl, above=2pt]{không: gửi lại} (8.2,-6.1) |- ($(w.east)+(0,0.2)$);
\end{tikzpicture}
\caption{Luồng gửi tin cậy, xử lý ack và gửi lại trong \texttt{ReliabilityLayer}}
\label{fig:B-guilai}
\end{figure}

% -------------------------------------------------------------------------------------
\subsection{Kênh, phân mảnh và kiểm soát nghẽn}
\begin{figure}[H]
\centering
\begin{tikzpicture}[x=1cm, y=1cm]
  \node[bxB, text width=3.2cm] (in)  at (3.2,0)   {Datagram PAYLOAD\\đã giải header};
  \node[bxB, text width=2.6cm] (in2) at (10.4,0)  {Datagram FRAGMENT};
  \node[bxB, text width=3.6cm] (fa)  at (10.4,-1.6) {FragmentAssembler\\{\footnotesize tối đa 8 nhóm, huỷ sau 2 s}};
  \node[dec, text width=2.0cm] (e)   at (6.8,-3.6) {channelId\\trong envelope};
  \node[bxG, text width=2.4cm] (c0)  at (1.8,-6.0) {Giao ngay\\không thứ tự};
  \node[dec, text width=2.3cm] (c1)  at (6.8,-6.0) {channelSequence\\mới hơn?};
  \node[dec, text width=2.3cm] (c2)  at (12.6,-6.0) {đúng sequence\\kế tiếp?};
  \node[bxG, text width=2.2cm] (d1)  at (5.2,-8.6) {Giao};
  \node[bxO, text width=2.2cm] (x1)  at (8.4,-8.6) {Huỷ gói cũ};
  \node[bxG, text width=2.8cm] (d2)  at (11.2,-8.6) {Giao + xả gói đệm liền sau};
  \node[bxO, text width=2.0cm] (b2)  at (14.5,-8.6) {Đệm chờ};
  \draw[arr] (in2) -- (fa);
  \draw[arr] (in.south) -- (3.2,-2.4) -| (e.north);
  \draw[arr] (fa.south) -- (10.4,-2.4) -| (e.north);
  \draw[arr] (e.west) -| node[lbl, pos=0.25, above=2pt]{kênh 0} (c0.north);
  \draw[arr] (e) -- node[lbl, right=2pt]{kênh 1 / 3} (c1);
  \draw[arr] (e.east) -| node[lbl, pos=0.25, above=2pt]{kênh 2} (c2.north);
  \draw[arr] (c1.south) -- ++(0,-0.4) -| node[lbl, pos=0.75, left=2pt]{có} (d1.north);
  \draw[arr] (c1.south) -- ++(0,-0.4) -| node[lbl, pos=0.75, right=2pt]{không} (x1.north);
  \draw[arr] (c2.south) -- ++(0,-0.4) -| node[lbl, pos=0.75, left=2pt]{có} (d2.north);
  \draw[arr] (c2.south) -- ++(0,-0.4) -| node[lbl, pos=0.75, right=2pt]{sớm} (b2.north);
\end{tikzpicture}
\caption{Xử lý phía nhận theo kênh: kênh 2 chờ đúng thứ tự, kênh 1 và 3 bỏ gói cũ}
\label{fig:B-kenh}
\end{figure}

\begin{figure}[H]
\centering
\begin{tikzpicture}[x=1cm, y=1cm]
  \node[init] (i) at (0.6,0) {};
  \node[st, fill=cGreen, text width=2.2cm] (g) at (3.2,0) {GOOD};
  \node[st, fill=cOrange, text width=2.2cm] (b) at (11.2,0) {BAD};
  \draw[arr] (i) -- (g);
  \draw[arr] ($(g.east)+(0,0.2)$) -- node[lbl, above=3pt]{RTT làm mượt > 250 ms} ($(b.west)+(0,0.2)$);
  \draw[arr] ($(b.west)+(0,-0.2)$) -- node[lbl, below=3pt]{RTT < 200 ms và đã ở BAD $\ge$ 10 s} ($(g.east)+(0,-0.2)$);
  \node[font=\footnotesize, align=center, anchor=north] at ($(g.south)+(0,-0.25)$) {snapshot 20 Hz};
  \node[font=\footnotesize, align=center, anchor=north] at ($(b.south)+(0,-0.25)$) {snapshot 10 Hz, giảm chi tiết};
\end{tikzpicture}
\caption{Kiểm soát nghẽn hai chế độ có trễ (hysteresis) 250/200 ms; RTT làm mượt theo EWMA 0,9/0,1}
\label{fig:B-nghen}
\end{figure}

% -------------------------------------------------------------------------------------
\subsection{Giao diện lập trình và giao diện chẩn đoán}
\begin{figure}[H]
\centering
\begin{tikzpicture}[x=1cm, y=1cm,
  cls/.style={draw=cLine, rectangle split, rectangle split parts=2, rectangle split part fill={cBlue,white},
    font=\footnotesize, inner sep=4pt}]
  \node[cls, anchor=north] (ic) at (2.9,0) {\begin{tabular}{@{}c@{}}\textit{«interface»}\\\textbf{ITransportClient}\end{tabular}
    \nodepart{two}\begin{tabular}{@{}l@{}}
      + State : ConnectionState \\ + SmoothedRttMs : float \\ + Stats : TransportStats \\
      + Connect(host, port, joinTicket) \\ + Send(channelId, payload, reliable) \\ + Disconnect() \\ + Poll() \\
      + OnMessage, OnConnected, \\ \quad OnDisconnected \\
    \end{tabular}};
  \node[cls, anchor=north] (is) at (8.4,0) {\begin{tabular}{@{}c@{}}\textit{«interface»}\\\textbf{ITransportServer}\end{tabular}
    \nodepart{two}\begin{tabular}{@{}l@{}}
      + ConnectionCount : int \\ + Start(port, maxConnections) \\ + Send(connectionId, channel, \\ \quad payload, reliable) \\
      + Broadcast(channel, payload, reliable) \\ + Disconnect(connectionId, reason) \\ + Poll() \\
      + OnValidateTicket, OnClientConnected, \\ \quad OnClientDisconnected \\
    \end{tabular}};
  \node[cls, anchor=north] (ts) at (13.9,0) {\textbf{TransportStats}
    \nodepart{two}\begin{tabular}{@{}l@{}}
      + BytesSent, BytesReceived \\ + PacketsSent, Received \\ + PacketsLost, Resent \\ + SmoothedRttMs, JitterMs \\
      + CongestionMode \\ + PendingFragmentGroups \\
    \end{tabular}};
  \node[bxY, text width=2.9cm] (uc) at (1.3,-7.4) {UdpTransportClient};
  \node[bxY, text width=2.9cm] (lb) at (4.6,-7.4) {LoopbackTransport};
  \node[bxY, text width=2.9cm] (us) at (8.4,-7.4) {UdpTransportServer};
  \draw[darr, -{Triangle[open, length=3mm]}] (uc.north) -- (uc.north |- ic.south);
  \draw[darr, -{Triangle[open, length=3mm]}] (lb.north) -- (lb.north |- ic.south);
  \draw[darr, -{Triangle[open, length=3mm]}] (us.north) -- (us.north |- is.south);
  \draw[arr] (is.east |- ts.west) -- (ts.west);
\end{tikzpicture}
\caption{Giao diện lập trình công khai của Transport (đóng băng từ tuần 1)}
\label{fig:B-api}
\end{figure}

\begin{figure}[H]
\centering
\begin{minipage}[t]{0.42\linewidth}
\begin{lstlisting}[basicstyle=\ttfamily\footnotesize]
Transport Connected
RTT 48.2 ms  jitter 3.1 ms
loss up 0.4%  down 1.2%
rate up 0.9 KB/s  down 1.7 KB/s
congestion GOOD  fragments 0  pool 12
packets sent 18422 recv 12291
        lost 147 resent 52
\end{lstlisting}
\end{minipage}\hfill
\begin{minipage}[t]{0.55\linewidth}
\begin{lstlisting}[basicstyle=\ttfamily\footnotesize]
$ dotnet run --project Ironfront.Tools.PacketReplay \
    -- capture.ifpcap --conn 3 --from 12.0 --to 15.0
t=12.004 OUT PAYLOAD  seq=4412 ack=3380 ch=1 len=87
t=12.021 IN  PAYLOAD  seq=3381 ack=4411 ch=3 len=41
t=12.062 IN  KEEPALIVE seq=3382 ack=4413 rtt~41ms
gap: seq 4414 not acked after 3 datagrams
summary: 118 out, 94 in, 2 gaps, rtt p50 44ms
\end{lstlisting}
\end{minipage}
\caption{Giao diện chẩn đoán: văn bản overlay do \texttt{TransportDiagnosticsFormatter} sinh (trái) và minh hoạ đầu ra
công cụ phát lại gói (phải); các số là minh hoạ bố cục}
\label{fig:B-chandoan}
\end{figure}

% =====================================================================================
\section{Luồng xử lý và giải pháp đã ứng dụng}
Mục này đi theo từng đầu việc của phân hệ trong bản mô tả đề tài: thiết lập kết nối UDP, header 16 byte, bắt
tay có cookie, ack bằng bitfield và gửi lại, bốn kênh truyền, cắt ghép mảnh, kiểm soát nghẽn, giữ kết nối, bộ giả
lập mạng và công cụ đo. Với mỗi đầu việc, báo cáo trình bày \textbf{luồng xử lý} (sơ đồ tuần tự theo đúng tên lớp
trong mã nguồn), \textbf{vấn đề gặp phải} và \textbf{giải pháp đã áp dụng}. Các tham số lấy trực tiếp từ mã nguồn
\code{Ironfront.Net.Transport} và \code{ProtocolConstants}.

% -------------------------------------------------------------------------------------
\subsection{Luồng thiết lập kết nối UDP có xác thực}
\textbf{Bài toán.} UDP không có kết nối: bất kỳ ai cũng gửi được datagram tới cổng 27015 và điền địa chỉ nguồn giả.
Nếu server cấp bộ nhớ, sinh số ngẫu nhiên hoặc kiểm chữ ký ngay khi nhận \code{CONNECT_REQUEST}, kẻ tấn công chỉ cần
gửi hàng loạt gói với địa chỉ nguồn giả là làm đầy bảng chờ và đẩy người chơi thật ra ngoài. Server cũng phải biết
người kết nối chính là người đã được Master Server cấp vé.

\begin{figure}[H]
\centering
\renewcommand{\seqstep}{0.66cm}
\begin{tikzpicture}[x=1cm, y=1cm]
  \seqhuman{B}{0.6}{NetClient-\\Bootstrap}
  \seqactor{C}{3.5}{c : Connection}
  \seqactor{S}{7.9}{srv : UdpTransport-\\Server}
  \seqactor{R}{11.1}{RateLimiter}
  \seqactor{K}{13.6}{Handshake-\\Cookie}
  \seqactor{V}{15.9}{Ticket-\\Validator}
  \seqbegin
  \seqmsg{B}{C}{BeginConnect(joinTicket)}
  \seqself{C}{clientSalt = CSPRNG u64\\trạng thái Connecting}
  \seqfragbegin
  \seqmsg{C}{S}{CONNECT\_REQUEST \{version, ticket, clientSalt\}}
  \seqfragend{C}{S}{loop}{mỗi 250 ms, tối đa 20 lần}
  \seqmsg{S}{R}{Allow(ip)}
  \seqret{R}{S}{true (5 yêu cầu/giây/IP)}
  \seqmsg{S}{K}{Derive(ip, port, clientSalt, epoch)}
  \seqret{K}{S}{serverSalt = HMAC-SHA256(...)[0..8]}
  \seqret{S}{C}{CONNECT\_CHALLENGE \{serverSalt\}}
  \seqself{C}{response = clientSalt XOR serverSalt\\trạng thái Challenged}
  \seqmsg{C}{S}{CONNECT\_RESPONSE \{response, clientSalt, ticket\}}
  \seqmsg{S}{K}{Verify(ip, port, clientSalt, response)}
  \seqret{K}{S}{true: địa chỉ nguồn đã được chứng minh}
  \seqmsg{S}{V}{OnValidateTicket(ticket)}
  \seqret{V}{S}{HMAC đúng, còn hạn, đúng serverId}
  \seqfragbegin
  \seqself{S}{cấp connectionId, gắn playerId,\\tạo Connection, lưu bản ghi phát lại}
  \seqret{S}{C}{CONNECT\_ACCEPTED \{connectionId, serverTick, mapId, playerId\}}
  \seqret{C}{B}{OnConnected(ConnectResult)}
  \seqfragelse{B}{V}{vé sai, hết hạn hoặc playerId đang chơi}
  \seqret{S}{C}{CONNECT\_DENIED \{InvalidTicket / AlreadyConnected\}}
  \seqret{C}{B}{OnDisconnected(lý do)}
  \seqfragend{B}{V}{alt}{hợp lệ}
  \seqend{B,C,S,R,K,V}
\end{tikzpicture}
\caption{Biểu đồ tuần tự thiết lập kết nối: server không lưu trạng thái cho tới khi địa chỉ nguồn được chứng minh}
\label{fig:B-luong-ketnoi}
\end{figure}

\textbf{Diễn giải luồng.}
\begin{enumerate}
  \item Client tạo \code{clientSalt} ngẫu nhiên bằng \code{RandomNumberGenerator} và gửi \code{CONNECT_REQUEST}
        (kích thước cố định, sai kích thước bị bỏ im lặng). Gói được gửi lại mỗi 250 ms, tối đa 20 lần
        (5 giây), vì chính gói này cũng có thể mất.
  \item Server kiểm tra theo thứ tự rẻ trước: giới hạn tốc độ theo IP, phiên bản giao thức, địa chỉ đã có kết nối
        chưa, còn chỗ trống không. Sau đó nó \textbf{không lưu gì} mà tính \code{serverSalt} bằng
        HMAC-SHA256 trên (địa chỉ, cổng, clientSalt, khung thời gian 30 s) và trả \code{CONNECT_CHALLENGE} về đúng
        địa chỉ đã khai.
  \item Chỉ client thật sự sở hữu địa chỉ đó mới nhận được challenge và tính được
        \code{clientSalt} XOR \code{serverSalt}. Khi \code{CONNECT_RESPONSE} tới, server tính lại cookie (chấp nhận khung
        hiện tại và khung liền trước) để xác minh --- tới đây địa chỉ nguồn mới đáng tin.
  \item Lúc này server mới làm việc tốn kém: kiểm HMAC của joinTicket, đọc \code{playerId} và từ chối nếu
        \code{playerId} đó đang có kết nối khác. Hợp lệ thì lấy một \code{connectionId} từ hàng đợi id trống,
        tạo \code{Connection} và trả \code{CONNECT_ACCEPTED}.
  \item \code{CONNECT_ACCEPTED} có thể mất. Server giữ một bản ghi nhỏ theo địa chỉ; nếu client gửi lại
        \code{CONNECT_RESPONSE} với đúng đáp án, server gửi lại \code{CONNECT_ACCEPTED} cũ thay vì cấp kết nối thứ hai.
\end{enumerate}

\begin{center}\small
\begin{tabularx}{\linewidth}{|L{4.4cm}|Y|L{3.6cm}|}
\hline
\hd{Vấn đề} & \hd{Giải pháp đã áp dụng} & \hd{Mã nguồn} \\ \hline
Lụt \code{CONNECT_REQUEST} với địa chỉ nguồn giả làm đầy bảng chờ & Cookie không trạng thái kiểu SYN-cookie: \code{serverSalt} là HMAC của địa chỉ, cổng, clientSalt và epoch 30 s, tính lại khi cần, không bao giờ lưu & \texttt{HandshakeCookie.Derive / Verify} \\ \hline
Server bị dùng làm công cụ khuếch đại & Challenge chỉ gửi một gói nhỏ về địa chỉ đã khai; giới hạn 5 yêu cầu/giây/IP, bảng giới hạn tối đa 10\,000 IP và tự dọn sau 10 s & \code{RateLimiter} \\ \hline
Một vé bị lộ mở ra nhiều người chơi & Gắn \code{playerId} trong vé với \code{connectionId}; vé thứ hai bị từ chối \code{AlreadyConnected}, không đá người đang chơi & \code{_playerIdToConnection} \\ \hline
Mất \code{CONNECT_ACCEPTED} & Bản ghi phát lại theo địa chỉ, hết hạn sau 10 s & \code{AcceptedHandshake} \\ \hline
Con trỏ ack bị gieo sai từ bộ đếm gói điều khiển dùng chung & Gói bắt tay \textbf{không} cập nhật con trỏ ack; gói dữ liệu đầu tiên mới khởi tạo nó & \code{Connection.HandleClientHandshake} \\ \hline
\end{tabularx}
\end{center}

% -------------------------------------------------------------------------------------
\subsection{Luồng gửi và nhận một bản tin}
Tầng trên (Replication) chỉ gọi \texttt{Send(channelId, payload, reliable)} và nhận lại \code{OnMessage(payload)}.
Mọi việc ở giữa --- đánh số, dựng header, ghi nhận để gửi lại, đưa vào socket, phân loại theo kênh --- là của
Transport. Toàn bộ chạy trên luồng gọi \code{Poll()}, không tạo luồng riêng, nên không cần khoá.

\begin{figure}[H]
\centering
\renewcommand{\seqstep}{0.66cm}
\begin{tikzpicture}[x=1cm, y=1cm]
  \seqhuman{A}{0.6}{Replication\\(bên gửi)}
  \seqactor{C}{3.4}{c : Connection}
  \seqactor{L}{6.6}{ReliabilityLayer}
  \seqactor{P}{9.7}{UdpPeer}
  \seqactor{Q}{12.6}{peer : Connection\\(bên nhận)}
  \seqactor{H}{15.5}{ChannelSet}
  \seqbegin
  \seqmsg{A}{C}{Send(kênh 2, payload, reliable)}
  \seqself{C}{kiểm CanSendReliable (< 64 gói chưa ack)\\ghi envelope \{channelId, channelSequence\}}
  \seqmsg{C}{L}{BuildAck(), NextSequence()}
  \seqret{L}{C}{ack, ackBitfield, sequence}
  \seqmsg{C}{L}{OnPacketSent(seq, datagram, reliable)}
  \seqnote{L}{L}{}{chép datagram vào BufferPool,\\giữ tới khi được ack}
  \seqmsg{C}{P}{Send(datagram, endpoint)}
  \seqmsg{P}{Q}{SendTo() -- datagram UDP $\le$ 1200 byte}
  \seqnote{P}{Q}{}{TryParse header: protocolId, cờ dự trữ = 0, connectionId khớp địa chỉ; sai thì bỏ im lặng}
  \seqselfl{Q}{ProcessIncomingAck(ack, bitfield)\\OnPacketReceived(sequence)}
  \seqmsg{Q}{H}{Receive(channelId, channelSequence, body)}
  \seqret{H}{Q}{giao theo chính sách kênh}
  \seqret{Q}{A}{KEEPALIVE mang ack mới (vì gói tin cậy)}
  \seqend{A,C,L,P,Q,H}
\end{tikzpicture}
\caption{Biểu đồ tuần tự gửi -- nhận một bản tin tin cậy qua kênh 2}
\label{fig:B-luong-guinhan}
\end{figure}

\textbf{Các giải pháp trên đường nóng.}
\begin{itemize}
  \item \textbf{Ack cõng trên mọi gói (piggyback).} Header nào cũng mang \code{ack} và \code{ackBitfield} 32 bit,
        nên một gói xác nhận được 33 gói cùng lúc và không cần gói ACK riêng. Riêng khi nhận gói tin cậy, bên nhận
        gửi ngay một \code{KEEPALIVE} để bên gửi không phải chờ tới nhịp giữ kết nối 1 giây.
  \item \textbf{Không cấp phát trên đường nóng.} Mọi mảng byte lấy từ \code{BufferPool} kích thước cố định 1200 byte.
        Đo trên .NET 8 Release: \code{BufferPool} 10,4 ns/lần, \code{ArrayPool.Shared} 17,4 ns, \texttt{new byte[]}
        16,7 ns; sau khởi động, \code{Poll} và \code{Send} cấp phát 0 byte.
  \item \textbf{Nhận có giới hạn.} Mỗi lần \code{Poll} đọc tối đa 1024 gói; phần còn lại để lượt sau (33 ms).
        Không giới hạn thì khi bị lụt gói, vòng mô phỏng nằm trọn trong vòng nhận và server ngừng mô phỏng.
  \item \textbf{Bộ đệm socket 1 MB và cảnh báo khi hệ điều hành cắt bớt.} Linux mặc định giới hạn
        \code{rmem_max} khoảng 208 KB mà không báo lỗi; \code{UdpPeer} đọc lại kích thước thực tế và ghi cảnh báo.
        Trên Windows, tắt \code{SIO_UDP_CONNRESET} để gói ICMP ``port unreachable'' không làm hỏng lần nhận sau.
  \item \textbf{Gói lỗi không được trả lời.} Header sai, \code{connectionId} không khớp địa chỉ hoặc cờ dự trữ khác 0
        đều bị bỏ im lặng --- trả lời sẽ biến server thành công cụ dò cổng.
\end{itemize}

% -------------------------------------------------------------------------------------
\subsection{Luồng phát hiện mất gói và gửi lại}
\begin{figure}[H]
\centering
\renewcommand{\seqstep}{0.66cm}
\begin{tikzpicture}[x=1cm, y=1cm]
  \seqactor{L}{1.6}{sender : ReliabilityLayer}
  \seqactor{P}{6.4}{UdpPeer}
  \seqactor{R}{10.6}{receiver : Connection}
  \seqactor{H}{14.6}{ChannelSet}
  \seqbegin
  \seqmsg{P}{R}{seq 100 (kênh 2, channelSeq 7)}
  \seqmsg{R}{H}{giao channelSeq 7}
  \seqlost{P}{R}{seq 101 (kênh 2, channelSeq 8) -- MẤT}
  \seqmsg{P}{R}{seq 102 (kênh 2, channelSeq 9)}
  \seqmsg{R}{H}{channelSeq 9 tới sớm}
  \seqnote{H}{H}{}{đệm chờ channelSeq 8,\\chưa giao 9}
  \seqret{R}{P}{gói kế tiếp mang ack = 102, bitfield ...101}
  \seqmsg{P}{L}{ProcessIncomingAck(102, ...101)}
  \seqself{L}{đánh dấu 100, 102 đã nhận;\\mẫu RTT chỉ từ gói gửi lần đầu (Karn)}
  \seqfragbegin
  \seqself{L}{Update(): seq 101 quá RTO\\RTO = clamp(1,5 $\cdot$ SRTT + 4 $\cdot$ jitter, 30, 1000 ms)}
  \seqmsg{L}{P}{gửi lại datagram 101 từ BufferPool}
  \seqfragend{L}{P}{loop}{chưa ack và ResendCount $\le$ 10}
  \seqmsg{P}{R}{seq 101 (bản gửi lại, channelSeq 8)}
  \seqmsg{R}{H}{channelSeq 8}
  \seqret{H}{R}{giao 8 rồi xả 9 đang đệm}
  \seqend{L,P,R,H}
\end{tikzpicture}
\caption{Biểu đồ tuần tự mất gói trên kênh tin cậy: bitfield chỉ ra gói mất, RTO kích hoạt gửi lại, kênh giữ đúng thứ tự}
\label{fig:B-luong-mat}
\end{figure}

\begin{center}\small
\begin{tabularx}{\linewidth}{|L{4.2cm}|Y|}
\hline
\hd{Vấn đề} & \hd{Giải pháp đã áp dụng} \\ \hline
Biết gói nào mất mà không tốn gói ACK riêng & Cửa sổ nhận 32 bit dịch theo khoảng cách sequence; bên gửi quét \code{ack} và 32 bit để xác nhận từng gói (\code{ProcessIncomingAck}) \\ \hline
Số thứ tự 16 bit quay vòng sau khoảng 36 phút ở 30 gói/giây & So sánh bằng \texttt{SequenceMath.IsNewer / Distance} theo số học modulo $2^{16}$, không dùng phép so sánh thường; có test biên tại 65535 $\rightarrow$ 0 \\ \hline
RTT bị đo sai khi gói đã gửi lại & Thuật toán Karn: chỉ lấy mẫu RTT khi gói được ack mà chưa từng gửi lại; SRTT làm mượt EWMA 0,9/0,1, jitter theo hệ số 0,25 \\ \hline
Chờ gửi lại quá lâu hoặc quá dồn & RTO động = 1,5 $\times$ SRTT + 4 $\times$ jitter, kẹp trong [30; 1000] ms \\ \hline
Gói tin cậy mất vĩnh viễn làm kênh có thứ tự treo mãi & Quá 10 lần gửi lại hoặc trùng ô vòng đệm 1024 thì \textbf{đóng kết nối} với \code{TransportError} thay vì tiếp tục im lặng (kết nối trông khoẻ nhưng không giao gì) \\ \hline
Bên nhận quá tải & Điều khiển luồng: tối đa 64 gói tin cậy chưa ack; \code{KEEPALIVE} quảng bá số gói chờ và mức áp lực, bên gửi tạm dừng gói tin cậy mới khi áp lực > 80\% \\ \hline
\end{tabularx}
\end{center}

% -------------------------------------------------------------------------------------
\subsection{Bốn kênh truyền và luồng phân loại khi nhận}
Mỗi gói \code{PAYLOAD} có envelope 3 byte gồm \code{channelId} và \code{channelSequence} \textbf{riêng từng kênh},
tách khỏi sequence của kết nối. Nhờ vậy \code{KEEPALIVE} hay gói của kênh khác không làm một snapshot trông ``cũ''.

\begin{center}\small
\begin{tabularx}{\linewidth}{|L{1cm}|L{3.8cm}|Y|L{3.8cm}|}
\hline
\hd{Kênh} & \hd{Đảm bảo} & \hd{Cách xử lý phía nhận (\code{ChannelSet.Receive})} & \hd{Dữ liệu} \\ \hline
0 & Không tin cậy, không thứ tự & Giao ngay & Ping/pong \\ \hline
1 & Không tin cậy, có thứ tự & Giao nếu \code{channelSequence} mới hơn gói đã giao; gói cũ bị huỷ và đếm \code{StalePacketsDropped} & Snapshot server $\rightarrow$ client \\ \hline
2 & Tin cậy, có thứ tự & Cửa sổ 256 ô: gói đúng thứ tự thì giao và xả các gói liền sau đang đệm; gói sớm được chép vào bộ đệm; gói trùng hoặc quá cửa sổ bị bỏ & Sự kiện, spawn, chat \\ \hline
3 & Không tin cậy, có thứ tự & Như kênh 1 nhưng bộ đếm riêng & Input client $\rightarrow$ server \\ \hline
\end{tabularx}
\end{center}
\textbf{Giải pháp chống head-of-line blocking:} độ tin cậy áp dụng \textbf{theo kênh}. Thí nghiệm
\code{--phase4-report} cho thấy khi kênh 2 đang thiếu sequence 0, snapshot kênh 1 vẫn được giao ngay (1 gói), còn
kênh 2 đệm gói sau và chỉ giao cả hai khi gói thiếu tới. Nếu dùng TCP, snapshot phải chờ sự kiện bị mất.

% -------------------------------------------------------------------------------------
\subsection{Luồng cắt và ghép mảnh bản tin lớn}
Datagram giới hạn 1200 byte để không bị phân mảnh ở tầng IP trên đường có PPPoE, VPN hay đường hầm (payload tối đa
1184 byte sau header 16 byte). Bản tin lớn hơn --- ví dụ danh sách người chơi khi mới vào trận --- được cắt ở tầng
Transport.

\begin{figure}[H]
\centering
\renewcommand{\seqstep}{0.66cm}
\begin{tikzpicture}[x=1cm, y=1cm]
  \seqhuman{A}{0.6}{Replication}
  \seqactor{C}{4.0}{c : Connection}
  \seqactor{Q}{9.0}{peer : Connection}
  \seqactor{F}{13.2}{FragmentAssembler}
  \seqbegin
  \seqmsg{A}{C}{Send(kênh 2, payload 3000 byte)}
  \seqself{C}{envelope + payload > 1184 byte:\\groupId mới, count = 3 ($\le$ 64)}
  \seqfragbegin
  \seqmsg{C}{Q}{FRAGMENT \{groupId, index i, count 3\} -- cờ RELIABLE | FRAGMENTED}
  \seqmsg{Q}{F}{TryReassemble(groupId, i, 3, data)}
  \seqret{F}{Q}{chưa đủ mảnh}
  \seqfragend{A}{F}{loop}{i = 0..2, mỗi mảnh một gói tin cậy}
  \seqfragbegin
  \seqret{F}{Q}{đủ 3 mảnh: ghép thành một mảng}
  \seqselfl{Q}{ProcessPayload: đọc envelope,\\đưa vào ChannelSet}
  \seqret{Q}{A}{OnMessage(payload 3000 byte)}
  \seqfragelse{A}{F}{nhóm quá 2 s hoặc đã có 8 nhóm dở dang}
  \seqselfl{F}{huỷ nhóm cũ nhất, trả mảnh về BufferPool}
  \seqfragend{A}{F}{alt}{đủ mảnh}
  \seqend{A,C,Q,F}
\end{tikzpicture}
\caption{Biểu đồ tuần tự cắt và ghép mảnh với bộ nhớ có giới hạn}
\label{fig:B-luong-manh}
\end{figure}

\textbf{Giải pháp.} (1) Mảnh luôn được gửi \textbf{tin cậy} kể cả khi tầng trên yêu cầu không tin cậy, vì thiếu một mảnh
thì cả bản tin vô dụng. (2) Giới hạn bộ nhớ trước dữ liệu do kẻ tấn công điều khiển: tối đa 64 mảnh mỗi nhóm, 8 nhóm
dở dang mỗi kết nối, nhóm quá 2 giây bị huỷ, và nếu \code{count} của một mảnh khác với nhóm thì huỷ cả nhóm.
(3) Trước khi cắt, kiểm tra số gói tin cậy chưa ack cộng số mảnh không vượt 64, tránh một bản tin lớn chiếm hết cửa sổ.
Kiểm thử: ghép đúng từng byte khi mảnh đến ngược thứ tự với bản tin 20 KB.

% -------------------------------------------------------------------------------------
\subsection{Luồng giữ kết nối, kiểm soát nghẽn và ngắt kết nối}
\begin{figure}[H]
\centering
\renewcommand{\seqstep}{0.66cm}
\begin{tikzpicture}[x=1cm, y=1cm]
  \seqhuman{T}{0.6}{Vòng tick\\(Poll)}
  \seqactor{C}{4.2}{c : Connection}
  \seqactor{G}{8.2}{CongestionControl}
  \seqactor{F}{11.9}{FlowControl}
  \seqactor{Q}{15.0}{peer}
  \seqbegin
  \seqfragbegin
  \seqmsg{T}{C}{Update(nowMs)}
  \seqmsg{C}{G}{Update(dt, SmoothedRttMs)}
  \seqret{G}{C}{GOOD: 20 Hz / BAD: 10 Hz, giảm chi tiết}
  \seqfragbegin
  \seqmsg{C}{Q}{KEEPALIVE \{pendingReliable u16, pressure \%\}}
  \seqfragend{C}{Q}{opt}{$\ge$ 1 s từ keep-alive trước}
  \seqret{Q}{C}{KEEPALIVE của peer}
  \seqmsg{C}{F}{ApplyRemote(pending, pressure)}
  \seqret{F}{C}{tạm dừng gói tin cậy mới nếu pressure > 80\%}
  \seqfragend{T}{Q}{loop}{mỗi tick khi Connected}
  \seqfragbegin
  \seqself{C}{Fail(Timeout): trả buffer, xoá kênh, xoá mảnh}
  \seqret{C}{T}{OnDisconnected(Timeout)}
  \seqfragelse{T}{Q}{Disconnect() chủ động}
  \seqmsg{C}{Q}{DISCONNECT \{lý do\} $\times$ 3 lần}
  \seqret{C}{T}{OnDisconnected(LocalRequest)}
  \seqfragelse{T}{Q}{gói tin cậy bị bỏ sau 10 lần gửi lại}
  \seqret{C}{T}{OnDisconnected(TransportError)}
  \seqfragend{T}{Q}{alt}{không nhận gì trong 10 s}
  \seqend{T,C,G,F,Q}
\end{tikzpicture}
\caption{Biểu đồ tuần tự giữ kết nối, điều khiển luồng và các đường kết thúc kết nối}
\label{fig:B-luong-giu}
\end{figure}

\begin{center}\small
\begin{tabularx}{\linewidth}{|L{4.2cm}|Y|}
\hline
\hd{Vấn đề} & \hd{Giải pháp đã áp dụng} \\ \hline
NAT xoá ánh xạ cổng khi kết nối im lặng; phát hiện client đã mất & \code{KEEPALIVE} mỗi 1 s; không nhận gói nào trong 10 s thì coi như rời trận (đúng yêu cầu ``mất kết nối quá 10 giây'' trong đề) \\ \hline
Keep-alive không bao giờ được gửi trên kết nối đang bận & Hẹn giờ theo lần gửi \textbf{keep-alive} cuối, không theo lần gửi gói bất kỳ; nếu không, kết nối gửi snapshot 20 Hz không bao giờ gửi keep-alive và thông tin điều khiển luồng không được làm mới \\ \hline
Báo ngắt kết nối có thể mất & \code{DISCONNECT} gửi 3 lần liên tiếp, không cần ack; bên kia vẫn có lưới an toàn timeout 10 s \\ \hline
Mạng nghẽn làm RTT tăng & Hai chế độ có trễ (hysteresis): SRTT > 250 ms chuyển BAD (snapshot 10 Hz, giảm chi tiết); về GOOD khi SRTT < 200 ms và đã ở BAD đủ 10 s; nếu vừa về GOOD chưa tới 10 s đã xấu lại thì phạt 20 s để không dao động \\ \hline
Rò rỉ tài nguyên khi đóng & \code{Fail} idempotent: trả mọi buffer của \code{ReliabilityLayer}, \code{ChannelSet}, \code{FragmentAssembler}; server trả \code{connectionId} về hàng đợi và gỡ ràng buộc \code{playerId} để người chơi kết nối lại được \\ \hline
\end{tabularx}
\end{center}

% -------------------------------------------------------------------------------------
\subsection{Bộ giả lập mạng, ghi phát lại gói và đo hiệu năng}
\textbf{Bộ giả lập mạng.} \code{NetworkSimulator} nằm ngay dưới \code{UdpPeer.Send}: mỗi datagram được tung xúc xắc
để bỏ, nhân đôi hoặc giữ lại một khoảng trễ (có jitter, có thể bị đẩy chậm thêm để đảo thứ tự), rồi được xả theo
thời điểm đến hạn trong \code{Poll}. Hạt giống cố định nên cùng cấu hình cho cùng chuỗi sự cố --- lỗi tái hiện được.
Bật mà không cần build lại: \code{IRONFRONT_SIM=typical}.

\begin{center}\small
\begin{tabular}{|l|c|c|c|c|c|}
\hline
\hd{Hồ sơ} & \hd{Trễ} & \hd{Jitter} & \hd{Mất gói} & \hd{Đảo thứ tự} & \hd{Trùng gói} \\ \hline
\code{lan} & 1 ms & 0 & 0 & 0 & 0 \\ \hline
\code{good} & 30 ms & 5 ms & 0,5\% & 0 & 0 \\ \hline
\code{typical} & 50 ms & 20 ms & 5\% & 2\% & 0 \\ \hline
\code{bad} & 100 ms & 50 ms & 15\% & 5\% & 2\% \\ \hline
\code{awful} & 150 ms & 100 ms & 30\% & 10\% & 5\% \\ \hline
\end{tabular}
\end{center}

\textbf{Ghi và phát lại gói.} Đặt \code{IRONFRONT_PCAP=đường-dẫn} thì \code{PacketLogger} ghi mọi datagram vào tệp
\code{.ifpcap} có phiên bản (chiều, thời điểm, địa chỉ, byte thô). Công cụ \code{Ironfront.Tools.PacketReplay} đọc
lại, lọc theo khoảng thời gian và \code{connectionId}, giải header, ước lượng lỗ hổng sequence và ghép ack với RTT
(Hình~\ref{fig:B-chandoan}). Vì tệp có thể chứa joinTicket, tệp chỉ lưu trong thư mục bị git bỏ qua.

\textbf{Kết quả đo (localhost, .NET 8 Release, ghi trong \code{plans/transport/reports/measurements.csv}).}
\begin{center}\small
\begin{tabularx}{\linewidth}{|L{5.6cm}|Y|}
\hline
\hd{Phép đo} & \hd{Kết quả} \\ \hline
Phân tích header GSP & 3,5 ns/lần, 0 byte cấp phát \\ \hline
16 kết nối, 30 Hz, 10 phút & 16/16 kết nối ổn định, bộ đệm không phải cấp thêm \\ \hline
64 kết nối, 30 Hz, 20 s (đã khởi động) & 64/64 kết nối; CPU 0,39\% một nhân; \code{Send} và \code{Poll} cấp phát 0 byte \\ \hline
Ack bitfield bật / tắt & Một gói ack xác nhận 33 gói / 1 gói; còn 7 / 39 gói tin cậy chờ \\ \hline
Kiểm soát nghẽn & RTT cao: 20 $\rightarrow$ 10 Hz; thời gian ở BAD tối thiểu 20 s khi vừa chuyển trạng thái \\ \hline
\end{tabularx}
\end{center}

\subsection{Kiểm thử phân hệ}
Bộ \code{Ironfront.Net.Transport.Tests} gồm 85 test xUnit, chạy trong CI trên Ubuntu và Windows:
\begin{center}\small
\begin{tabularx}{\linewidth}{|L{4.6cm}|c|Y|}
\hline
\hd{Tệp test} & \hd{Số test} & \hd{Kịch bản tiêu biểu} \\ \hline
\code{ControlAndSocketTests} & 22 & Gói điều khiển, timeout 10 s, DISCONNECT, socket thật trên localhost \\ \hline
\code{ReliabilityLayerTests} & 19 & Ack bitfield, sequence quay vòng, Karn, RTO, bỏ gói sau 10 lần \\ \hline
\code{HandshakeHardeningTests} & 12 & Cookie sai, cookie hết epoch, vé sai, vé trùng playerId, lụt yêu cầu \\ \hline
\code{ChannelAndFragmentTests} & 11 & Sai thứ tự, gói trùng, gói cũ trên kênh 1/3, ghép mảnh ngược thứ tự, giới hạn 8 nhóm \\ \hline
\code{ReliableChannelSurvivalTests} & 8 & Kênh tin cậy sống sót qua mất gói kéo dài của bộ giả lập \\ \hline
\code{UdpPeerTests} & 8 & Ngân sách 1024 gói/Poll, bộ đệm socket, gói header sai \\ \hline
\code{PacketLoggerTests}, \code{Diagnostics} & 5 & Ghi -- đọc lại \code{.ifpcap}, định dạng overlay chẩn đoán \\ \hline
\end{tabularx}
\end{center}
